\begin{frame}{Resumo da DP da moeda}
    \metroset{block=fill}

    Resumo do processo:
    \begin{itemize}
        \item Escolha estados que representem seu problema.
        \item Escreva uma função que resolva o problema a partir de subproblemas.
        \item Calcule a complexidade. Se a transição de cada estado tem a mesma complexidade, então a complexidade total é $O(estados * transicao)$.
    \end{itemize}

    No problema anterior, decidimos escolher as moedas da esquerda para a direita. Então o estado deve ser $E = \{0, 1, \dots, n\}$, onde $i \in E$ representa que ainda temos as moedas $i, i+1, \dots, n-1$ para escolher.

    A função recursiva que resolve o problema é $$\DP(i) = \max(\DP(i+1), v_i + \DP(i+2)) \ .$$.

    Como a transição de todo estado é constante, a complexidade é $O(n)$.
\end{frame}
